% ====================================================================
% Robustness of the tube/sp regression: leave-one-out + Cook's distance.
% Answers Kimi (feared the correlation leans on BA points; opposite is true).
% Goes into the tube/sp applicability section.
% Data: experiments/leverage_tubesp.py
% ====================================================================
\subsection{Is the predictor robust, or does it lean on a few points?}
\label{sec:leverage}

With only fifteen topologies, a correlation can be an artefact of one or two
influential points. We checked. Across all single-point deletions
(leave-one-out), the correlation coefficient moves by at most $0.059$; no
individual topology drives the result. Only one point, the densest scale-free
graph BA$_{n=50,m=3}$, exceeds the conventional Cook's distance cutoff of
$4/n$, and removing it \emph{raises} the correlation rather than lowering it.

We also ran the specific stress test a skeptic would ask for: removing all three
Barab\'asi--Albert points, which sit visibly below the regression line. Far from
collapsing, the correlation strengthens---$r$ rises from $0.78$ to $0.89$
($R^2$ from $0.60$ to $0.79$). The scale-free graphs were the weakest-fitting
points, slightly flattening the relationship; the structural predictor is in
fact cleaner on the grid and small-world families than the full-sample $r=0.78$
suggests. We report the full-sample figure as the conservative one. The point
stands: tube/sp predicts the room a topology leaves for routing to matter, and
the prediction does not hang on any handful of graphs.

% ====================================================================
% Out-of-sample validation on SNDlib backbones (bench A).
% Data: experiments/bench_a_v3.py (gated by positive control).
% ====================================================================
\subsection{Does the predictor hold out of sample?}
\label{sec:oos}

The fifteen bench topologies were also the ones used to read off the
relationship. A predictor is only worth the name if it survives on data it
never saw. We therefore took eighteen real backbone topologies from
SNDlib---national and continental research and carrier networks, none of them
in the bench---and asked whether tube/sp still predicts where the physical core
gains over CONGA.

To keep the test fair, every SNDlib graph receives the same capacity and
latency scheme as the bench, so the only thing that varies from the training
set is the structure itself, which is exactly the quantity under test. We also
guard the measurement with a positive control: before trusting any
out-of-sample number, the identical pipeline is run on the bench topologies and
must reproduce the known in-sample relationship. It does ($r=0.62$ on the
bench through this path), so the instrument is faithful.

On the eighteen unseen backbones the relationship is not merely preserved, it
is sharper than in sample: tube/sp and loss reduction correlate at Pearson
$r=0.85$ and Spearman $\rho=0.98$ ($p<10^{-4}$). The sign structure the law
predicts appears cleanly. The five most rigid graphs---those with the lowest
tube/sp, where there is little alternative-path room to exploit---show the
physical core \emph{losing} to CONGA, exactly as a structural ceiling on
routing gain implies. As tube/sp rises, the advantage climbs monotonically,
reaching $+63\%$ on the topology with the most routing room.

The one graph that sits off the line, dfn-gwin, has the highest tube/sp of all
yet only a middling gain; with eleven nodes its high ratio reflects a small
graph rather than genuine path richness---the same small-$n$ caveat that
motivated the ratio in the first place. We read the result as the predictor
doing what a structural predictor should: ordering topologies by how much room
they leave for routing to matter, on networks it was never fitted to.
